% Auto-generated by experiments_v2/process_duke_results.py
% Replaces \needdata{} placeholders in tab:duke-summary and
% tab:duke-probe-vs-l1.  Re-run after every cluster sweep.

\begin{table}[h]
  \centering
  \caption{Achievable selectivity index $\mathrm{SI}$ on the Duke microCT-derived vagus nerve cohort, by species and optimisation path.  ``Probe'' is the probe-based smart-init Adam-FD; ``L1'' is the random-init L1-discovery validation path.  Cells: mean / median / \% reaching SI $\geq 0.95$.}
  \label{tab:duke-summary}
  \begin{tabular}{lrlll}
    \toprule
    Species & $n$ & Probe SI & L1 SI & Wave SI \\
    \midrule
    swine & 4 & 0.890 / 0.896 / 0\% & 0.619 / 0.749 / 25\% & 0.891 / 0.896 / 0\% \\
    human & 5 & 0.927 / 0.977 / 80\% & 0.808 / 0.876 / 40\% & 0.726 / 0.954 / 60\% \\
    all & 9 & 0.911 / 0.942 / 44\% & 0.724 / 0.862 / 33\% & 0.799 / 0.898 / 33\% \\
    \bottomrule
  \end{tabular}
\end{table}

\begin{table}[h]
  \centering
  \caption{Probe-based smart init vs random-init L1-discovery agreement on the Duke cohort.  Two paths are said to \emph{agree} if $|\Delta \mathrm{SI}| \leq 0.05$ (experimental noise floor).  $n_{\mathrm{active}}$ is the median number of non-zero contacts in the final amplitude vector.}
  \label{tab:duke-probe-vs-l1}
  \begin{tabular}{lrrrr}
    \toprule
    & mean $|\Delta \mathrm{SI}|$ & \% agree & probe $n_{\mathrm{active}}$ & L1 $n_{\mathrm{active}}$ \\
    \midrule
    swine & 0.280 & 50\% & 3 & 12 \\
    human & 0.120 & 40\% & 3 & 12 \\
    all & 0.191 & 44\% & 3 & 12 \\
    \bottomrule
  \end{tabular}
\end{table}
